\subsection{Sites, truth and label}
Six windows of USGS 3DEP one metre bare earth lidar~\cite{sugarbaker2014} cover active dune fields (Table~\ref{tab:sites}); we chose them from 20\,m slope images and had used them in earlier tests of the detector. A post is \emph{steep} when its Horn slope~\cite{horn1981} at 1\,m is $28^\circ$ or more, a little below the angle of repose of dry sand~\cite{pye2009}; voids never count as steep. A \emph{face patch} is an 8 connected group of at least 10 steep posts (10\,m$^2$). Face patches are steep ground, not verified slip faces. For a coarse cell $c$ with footprint $F_c$ the label is the face map
\begin{equation}
y_c = \begin{cases}1 & \text{if some } x \in F_c \text{ lies in a face patch},\\ 0 & \text{otherwise,}\end{cases}
\label{eq:label}
\end{equation}
which a planner needs and whose meaning does not change with the cell size. The cell label of the first version of this benchmark, at least 10 steep posts in $F_c$, means 10\,\% of a 10\,m cell but 1.1\,\% of a 30\,m cell and is reported only for comparison.

\subsection{Inputs}
\textit{Idealised:} block means of the lidar at 2, 5, 10, 15, 20, 30 and 60\,m (I10 and I30 in the tables) (no product noise, time gap or misregistration). \textit{Real:} the USGS 1/3 arc second DEM (R10, about $8\times10$\,m)~\cite{usgs3dep13}, NASADEM (N30, February 2000)~\cite{nasadem} and Copernicus GLO-30 (G30, 2011 to 2015)~\cite{copernicusdem}, each scored on its own geographic grid, with labels from the lidar posts inside each native cell; no horizontal shift is applied. For planning, GLO-30 is interpolated to the 30\,m lidar grid. Over these sites R10 is probably resampled by USGS from the same lidar collections. \textit{Imagery:} one Sentinel-2 L2A scene per site~\cite{drusch2012sentinel} (10\,m; lidar year, cloud below 5\,\%, lowest sun, chosen by rule) and NAIP~\cite{naip} (0.3 to 0.6\,m, resampled to 1\,m).

\subsection{Hazard layers}
Table~\ref{tab:methods} lists the layers, among them a logistic subgrid detector trained only on synthetic dune profiles, which ships with the planner and is one baseline among others. Every learned layer is trained leave one site out~\cite{roberts2017}.

\subsection{Planner}
The planner searches the 8 connected grid of the input. For an edge with along track grade $\tan\alpha$ (uphill positive) and cross slope $\tan\beta$,
\begin{equation}
d=\sqrt{(\tan\alpha+f_r)^2+\tan^2\beta},\qquad i=-i_0\ln\left(1-d/\tan\phi\right)
\label{eq:slip}
\end{equation}
give the traction demand and the slip ($i=\infty$ when $d\ge\tan\phi$), with friction angle $\phi$, motion resistance coefficient $f_r$ and slip scale $i_0$; the slip law is the exponential shear relation of terramechanics~\cite{wong2008} written for the ratio of demand to capacity. An edge is passable when $i\le i_{\max}$ and costs $\ell/(1-i)$ for length $\ell$. The four parameters are drawn 128 times; a cell is passable for a draw when at least four of its eight headings are, and $q_c$ is the share of draws for which it is. A cell is BLOCKED when $q_c\le0.05$, TRAVERSABLE when $q_c\ge0.95$ and its face probability $p_c<\tau$, and UNKNOWN otherwise. The main setting, the \emph{posterior vehicle}, draws the parameters from a posterior fitted by Markov chain Monte Carlo to published slip curves of other wheeled vehicles (5 to 95\,\% ranges: $\phi$ 30.6 to $43.7^\circ$, $f_r$ 0.03 to 0.32, $i_0$ 0.33 to 1.21; $i_{\max}$ 0.4 to 0.7; nominal climbing limit {{VAL:crit_posterior}}$^\circ$); a narrower assumed range (limit {{VAL:crit_assumed}}$^\circ$) is the sensitivity. Three policies share the graph: \emph{slip only}, \emph{gate} (TRAVERSABLE cells only) and \emph{weight} (no BLOCKED cell, cost times $1+w\,p_c$ of the cell entered); gate and weight refuse descents of $28^\circ$ or more. Every layer is run at $\tau\in\{0.02,\dots,0.9\}$ and $w\in\{1,\dots,100\}$ and reported at the setting that maximises mission success on the other five sites.

\subsection{Scoring and closed loop}
One hundred pairs per site, at least 1\,km apart, are drawn twice: E1 from cells whose lidar posts are all gentle (truth used, as in our earlier runs) and E2 from cells whose coarse slope is below $15^\circ$ (no truth). A route, the polyline through its cell centres sampled every 0.5\,m on the lidar, is \emph{unsafe} when a sample falls on a face patch, and a \emph{mission success} when it exists and is not unsafe. In the closed loop the rover follows the plan on a 5\,m grid that inherits the coarse classes and face probabilities ($w=5$), with slip costs from the coarse surface interpolated to 5\,m; at every step it senses the face patches within a radius $R$ from the lidar, blocks those cells, drops the face penalty of sensed free cells and replans when its route is blocked, within a window 300\,m around the pair and up to 300 replans.

\subsection{Statistics and prediction log}
AUC intervals come from a spatial block bootstrap~\cite{kunsch1989} with 300\,m blocks; route outcomes carry two stage bootstrap intervals (sites, then pairs). Site level claims use an exact sign flip test over the six sites, Holm corrected~\cite{holm1979}; the smallest attainable two sided $p$ is 1/32. On one core of a cloud x86 machine, the slip assessment of a {{VAL:rt_area}}\,km$^2$ window takes {{VAL:rt_I30_assess}}\,s at 30\,m and {{VAL:rt_I10_assess}}\,s at 10\,m, and planning 100 routes on it {{VAL:rt_I30_routes}} and {{VAL:rt_I10_routes}}\,s. Every result in Tables III to V was registered with numeric predictions in a prediction log before it ran, but after the results of the first version on the same sites were known; the checks in Section~\ref{sec:posthoc} are post hoc. The predictions of stages 25 to 28 ({{VAL:log_hit}} hits, {{VAL:log_missed}} misses) and three amendments made before the affected results were known are in the supplement.
